\documentclass[11pt]{article}
\usepackage{mzqastyle}
\setrunningtitle{SEC \& EDINET Metrics Mapping \& Reconciliation}

\begin{document}

\mzqatitlepage
  {Data-layer reference}
  {SEC \& EDINET Metrics Mapping \& Reconciliation}
  {From raw filed XBRL to a standardized, reconciled fundamentals layer}
  {The fundamentals AI Analyst reads are not the numbers companies file. They are
   the output of a multi-stage pipeline that ingests raw XBRL facts from SEC EDGAR
   (US) and EDINET (Japan), \emph{maps} tens of thousands of heterogeneous filer
   concepts onto a single governed set of standardized line items, \emph{derives}
   the items filers omit, \emph{checks} them against accounting identities,
   \emph{computes} a library of financial metrics, and \emph{reconciles} every
   computed figure back to the exact source concepts and filings it came from.

   \vspace{8pt}
   This volume specifies that pipeline end to end: the raw-to-standardized concept
   mapping and its governance model, the graph-closure derivation and
   accounting-identity checks, the metric formula registry and its evaluation, and
   the reconciliation and validation layers that make every number auditable.

   \vspace{8pt}
   \textit{Companion volumes:} \emph{System Architecture \& Operations},
   \emph{Quantitative \& ML Models}, and \emph{qlib Integration}.}
  {Metrics Mapping \& Reconciliation \quad·\quad Generated 31 July 2026}

\mzqatoc
\clearpage

%======================================================================
\part*{Part I\quad The standardized layer}
\addcontentsline{toc}{section}{Part I\quad The standardized layer}

\section{Why a standardized layer exists}

Two companies reporting the identical economic fact --- say, cost of goods sold ---
will tag it with different XBRL concepts, under different taxonomies, with
different signs, scales, and statement positions, and one of them may not tag a
subtotal the other reports explicitly. A US filer uses \code{us-gaap} concepts; a
Japanese filer uses the EDINET/\code{jppfs} taxonomy in Japanese. Neither is
directly comparable, and neither is stable across years as taxonomies evolve.

The \code{xbrl\_sec} warehouse solves this with a \emph{standardized layer}: a
single, governed vocabulary of standardized line items (\code{revenue},
\code{cost\_of\_goods\_sold}, \code{total\_financial\_debt}, \dots) onto which every
filer's raw concepts are mapped, and against which a uniform library of metrics is
computed. AI Analyst reads only this layer; the application never touches raw
filings. This document specifies how the layer is built and made auditable.

\begin{keybox}[The invariant]
Every standardized value and every computed metric can be traced back to the exact
raw concept, filing, and fiscal period it derives from --- or is explicitly marked
as \emph{derived} (from an accounting identity or a formula) rather than filed. The
reconciliation layer (\S\ref{sec:recon}) records that trace for every figure, and
the validation layer (\S\ref{sec:validate}) audits coverage and integrity.
\end{keybox}

\section{Pipeline overview}
\label{sec:overview}

Per jurisdiction the pipeline runs as an ordered set of stages. The application is
read-only; these stages are the write path that maintains the shared warehouse.

\begin{figure}[h]
\centering
\begin{tikzpicture}[node distance=4mm and 6.5mm, font=\sffamily\footnotesize,
   st/.style={proc, minimum width=20mm, minimum height=9mm}]
  \node[data, minimum width=22mm] (edgar) {SEC EDGAR\\\scriptsize companyfacts · XBRL};
  \node[data, below=8mm of edgar, minimum width=22mm] (edinet) {EDINET\\\scriptsize XBRL packages};
  \node[st, right=9mm of $(edgar)!0.5!(edinet)$] (parse) {Raw parse\\\scriptsize \code{fact\_fundamentals\_*}};
  \node[st, right=of parse] (std) {Standardize\\\scriptsize \code{..\_std\_*}};
  \node[st, right=of std] (metrics) {Metrics\\\scriptsize \code{fact\_metrics\_*}};
  \node[st, right=of metrics] (recon) {Recon\\\scriptsize \code{..\_recon\_*}};
  \node[st, right=of recon, fill=tealfill, draw=steel] (val) {Validate};

  \draw[arr] (edgar) -| (parse); \draw[arr] (edinet) -| (parse);
  \draw[arr] (parse) -- (std); \draw[arr] (std) -- (metrics);
  \draw[arr] (metrics) -- (recon); \draw[arr] (recon) -- (val);

  \node[below=1mm of parse, font=\sffamily\scriptsize\itshape, text=muted, align=center] {raw facts};
  \node[below=1mm of std, font=\sffamily\scriptsize\itshape, text=muted, align=center] {mapping +\\graph closure};
  \node[below=1mm of metrics, font=\sffamily\scriptsize\itshape, text=muted, align=center] {formula\\registry};
  \node[below=1mm of recon, font=\sffamily\scriptsize\itshape, text=muted, align=center] {source\\trace};
\end{tikzpicture}
\caption{The end-to-end fundamentals pipeline. Each stage is idempotent and
incrementally re-runnable; downstream stages have freshness gates so an unchanged
input is skipped.}
\label{fig:pipeline}
\end{figure}

\begin{center}\small
\begin{tabular}{@{}>{\bfseries\color{navy}}l p{6.0cm} >{\ttfamily\footnotesize}p{5.0cm}@{}}
\toprule
\tabhead{Stage} & \tabhead{What it produces} & \normalfont\tabhead{Primary tables} \\
\midrule
Raw parse    & Filed facts as (concept, value, unit, context, period, statement position) & fact\_fundamentals\_\{us,jp\} \\
Standardize  & Standardized line items via governed mapping $+$ graph closure & fact\_fundamentals\_std\_\{us,jp\} \\
Metrics      & Derived line items (L1) and financial metrics (L2) from the formula registry & fact\_metrics\_\{us,jp\} \\
Recon        & Every metric re-expressed with its inputs and traced to source concepts/filings & fact\_metrics\_recon\_\{us,jp\} \\
Validate     & Coverage, identity, sparsity, and master-quality reports & (report; ref\_std\_identity\_violation) \\
\bottomrule
\end{tabular}
\end{center}

The schema itself is defined as a large set of sequential, idempotent SQL
migrations, each guarded so re-running the set against a current database is safe.
The reference tables the standardizer depends on --- the standardized line-item
catalogue (\code{ref\_standardized\_line\_items}), the rollup edges
(\code{ref\_std\_item\_edge}), the identity checks (\code{ref\_std\_identity\_check}),
and the governed concept mappings --- are loaded before any fundamentals run.

%======================================================================
\part*{Part II\quad Raw acquisition and parsing}
\addcontentsline{toc}{section}{Part II\quad Raw acquisition and parsing}

\section{US: SEC companyfacts}
\label{sec:us-raw}
US fundamentals originate from SEC EDGAR's structured \emph{companyfacts} (and the
XBRL filing packages for statement structure). Each raw fact carries its
\code{us-gaap} (or filer-extension) concept, value, unit, the filing it came from,
the fiscal period, and --- from the calculation/presentation linkbases --- its
position in the statement tree (parent, root, weight, presentation order). These
are written to \code{fact\_fundamentals\_us}.

\paragraph{Fiscal-period alignment.} The SEC \code{fy} tag is the fiscal year of
the \emph{filing}, not of the reported period. A prior quarter that reappears as a
comparative in a later 10-Q or 10-K inherits the later filing's \code{fy}. Because
the standardized key is $(\text{cik},\text{fiscal\_year},\text{fiscal\_period},
\text{line\_item})$, two genuinely different quarters could otherwise collide.
Annual rows are therefore re-dated from their \code{period\_end}; quarterly rows are
period-aligned to the fiscal year they close into, using the company's inferred
fiscal-year-end month --- so a September-quarter fact for an early-FYE filer lands
in the correct fiscal year rather than the filing's.

\section{Japan: EDINET XBRL}
\label{sec:jp-raw}
Japanese fundamentals originate from EDINET XBRL packages. The parser
(\code{parsers/edinet\_xbrl.py}) reads the instance document together with the
calculation and presentation linkbases:

\begin{itemize}
  \item \textbf{Concept normalization} --- an \code{href} fragment such as
        \code{jppfs\_cor\_Sales} is normalized to a
        \code{prefix/localname} concept id; non-numeric concepts (text blocks,
        abstracts, axes, domains, members, tables) are skipped.
  \item \textbf{Statement classification} --- the linkbase \code{role} maps to a
        statement type (balance sheet / income statement / cash flow / equity) via
        role-string heuristics.
  \item \textbf{Context and dimensions} --- each fact carries an XBRL \emph{context}
        (period and any explicit dimensional members). JP standardization is
        context-aware: a \code{dimension\_signature} distinguishes, for example, a
        consolidated total from a segment or a prior-period restatement, so the
        right context is chosen per line item.
  \item \textbf{Identity concepts} --- cover-page DEI concepts (company name in
        English/Japanese, security code) are captured for entity resolution.
\end{itemize}

Facts are written to \code{fact\_fundamentals\_jp} with their context id,
dimension signature, and statement structure, mirroring the US raw shape so the
standardizer can treat both jurisdictions uniformly downstream.

%======================================================================
\part*{Part III\quad The standardization layer}
\addcontentsline{toc}{section}{Part III\quad The standardization layer}

\section{The governed mapping model}
\label{sec:mapmodel}
The heart of standardization is a governed table,
\code{map\_concept\_to\_taxonomy\_versioned}, that assigns each raw concept to a
standardized \code{target\_variable} under a set of qualifying conditions. Each
mapping row carries:

\begin{center}\small
\begin{tabular}{@{}>{\ttfamily\color{navy}}l p{10.3cm}@{}}
\toprule
\normalfont\tabhead{Field} & \normalfont\tabhead{Meaning} \\
\midrule
target\_variable & The standardized line item this concept feeds (e.g.\ \code{cost\_of\_goods\_sold}). \code{UNMAPPED} rows are excluded. \\
tier             & Mapping tier / precedence class (tier-1 primary vs.\ tier-2 secondary). \\
multiplier       & Scale/sign multiplier applied to the raw value. \\
sign\_policy / normal\_balance & How the concept's sign is interpreted against the standardized item's normal balance. \\
mapping\_sector  & The sector scope the mapping applies to (\code{corp}, \code{bank\_financial}, \code{insurance}, \code{reit}, \dots, or universal). \\
gics\_* & Optional GICS sector / industry-group / industry / sub-industry constraints. \\
effective\_from/to\_year & The fiscal-year range over which the mapping is valid (taxonomy evolution). \\
accounting\_standard & Required accounting standard (e.g.\ US-GAAP, IFRS, JGAAP), or unconstrained. \\
taxonomy\_version & Required taxonomy version, or unconstrained (year-level partial match allowed). \\
aggregation\_type / priority & How multiple concepts mapping to the same target combine. \\
confidence       & Governance confidence attached to the mapping. \\
\bottomrule
\end{tabular}
\end{center}

Two sources feed the resolver: the governed versioned table above, and approved
\emph{company-period exceptions}
(\code{map\_concept\_to\_taxonomy\_exception}, \code{review\_status='approved'},
unexpired) that override the general mapping for a specific entity and fiscal-year
range. Both are indexed by concept id with slash/colon aliasing so
\code{prefix/local} and \code{prefix:local} forms both resolve.

\section{Mapping resolution}
\label{sec:resolve}
For each raw fact, the resolver (\code{versioned\_mapping.select\_versioned\_mapping})
selects the single best mapping by a deterministic specificity order. An approved
exception matching the entity, year, period, sector, accounting standard, and
taxonomy wins outright. Otherwise, among mappings valid for the fact's fiscal year:

\begin{enumerate}
  \item \textbf{Accounting-standard \& taxonomy gate.} A mapping is discarded if
        its required accounting standard or taxonomy version is incompatible with
        the fact. Compatible matches are scored (exact standard/taxonomy $>$
        same-year taxonomy $>$ unconstrained), and these scores enter the ranking.
  \item \textbf{Sector specificity.} Candidates are considered in the entity's
        \emph{sector order} --- most specific first, then universal. Corporate
        filers fall back \code{corp}\,$\to$\,universal; banks, insurers, REITs, and
        asset managers use their own scopes plus universal (an insurer, for
        instance, resolves \code{insurance}\,$\to$\,\code{non\_bank\_financial}\,$\to$\,universal).
        The first sector with any qualifying candidate wins.
  \item \textbf{GICS specificity.} Within that sector, a mapping's GICS constraints
        must match the entity; more specific GICS constraints (sub-industry $>$
        industry $>$ group $>$ sector) rank higher.
  \item \textbf{Tie-breaks.} Remaining ties break by accounting/taxonomy score,
        then jurisdiction match, then earliest effective year, then confidence,
        then mapping id.
\end{enumerate}

\begin{notebox}[Why sector scope matters]
Sector scope prevents category errors that silently corrupt derived figures. An
insurer's operating investment portfolio, for example, must map to
\code{investment\_securities}, not to \code{short\_term\_investments} --- otherwise
it would be netted out as ``excess cash'' in corporate net-debt derivation. The
resolver's sector hierarchy, plus sector-specific rollup edges (\S\ref{sec:closure}),
keep each filer's economics on the right standardized items.
\end{notebox}

\section{Sign, unit, and value canonicalization}
\label{sec:canon}
Once a mapping is chosen, the raw value is canonicalized to the standardized item's
convention:
\begin{equation}
  v_{\text{std}} \;=\;
  \begin{cases}
    -\,\lvert v_{\text{raw}}\cdot \text{multiplier}\rvert, & \text{line item is canonically negative,}\\[2pt]
    \phantom{-}\,v_{\text{raw}}\cdot \text{multiplier}, & \text{otherwise.}
  \end{cases}
\end{equation}
Canonically negative items --- cost of goods sold, R\&D, SG\&A, interest expense,
tax provision, capex, dividends paid, share repurchases --- are stored with a
consistent negative sign regardless of how the filer signed them, so downstream
formulas can add line items without sign guesswork. Unit handling ranks candidate
facts so a monetary line item prefers a currency-unit fact over a share-count fact
of the same concept. The standardized row also records the source concept id,
filing id, concept path, standardized concept path, mapping id, currency, and (for
JP) the context id and dimension signature --- the raw material for the recon
trace.

\section{Graph closure: deriving what filers omit}
\label{sec:closure}
Filers report different subsets of a statement. After the raw mapping pass
populates the \emph{filed} leaves, \code{std/graph\_closure.py} walks the
standardized-item hierarchy (\code{ref\_std\_item\_edge}, filtered to the entity's
sector scope and accounting standard) to derive missing values at the
$(\text{entity},\text{fiscal\_year},\text{fiscal\_period})$ grain. It runs three
passes.

\subsection{Top-down residuals (catch-all children)}
When a parent total is known and exactly one child is missing --- preferentially a
residual/``catch-all'' child (item class \code{catch\_all}, derivation policy
\code{residual}, or an \code{other\_}$\ast$ item) --- that child is derived as the
parent minus its known siblings:
\begin{equation}
  \text{child} \;=\; \pm\Big(\text{parent} - \!\!\sum_{\text{known siblings}}\!\! \text{sign}\cdot\text{child}\Big),
\end{equation}
with the edge sign applied. This pre-pass runs first so an authoritative parent
total fills a residual before a partial bottom-up subtotal can.

\subsection{Bottom-up subtotals}
Intermediate totals are then computed from their children. A parent is marked
\code{DERIVED\_BOTTOM\_UP} when \emph{all} children are known, or
\code{DERIVED\_PARTIAL} when at least $80\%$ are known (below that it is left
missing rather than fabricated from too little). Expense signs are handled once:
a subtractive edge uses the magnitude of the child and applies the edge sign, so
already-negative expense facts are not double-counted.

\subsection{Formula line items (L1)}
Registry-defined derived line items --- the ones that are \emph{always} computed
rather than filed --- are then evaluated from known values in dependency order
(\S\ref{sec:registry}); for example
\begin{align}
  \text{total\_financial\_debt} &= \text{short\_term\_debt}+\text{LTD\_current}+\text{long\_term\_debt}, \\
  \text{net\_debt} &= \text{total\_financial\_debt}-\text{cash}-\text{short\_term\_investments}, \\
  \text{gross\_profit} &= \text{revenue}-\lvert\text{cost\_of\_goods\_sold}\rvert .
\end{align}
Sector overrides apply here too --- for insurers, \code{net\_debt} drops the
short-term-investment offset. Derived items carry a \code{metric\_type} tag
(\code{RESIDUAL}, \code{DERIVED\_BOTTOM\_UP}, \code{DERIVED\_PARTIAL},
\code{DERIVED\_TOP\_DOWN}) so their provenance is never confused with a filed value.

\section{Accounting-identity checks}
\label{sec:identities}
Closure also \emph{audits} the statements against accounting identities and emits
violations. Two families of check run at the same entity$\times$period grain.

\paragraph{Rollup consistency.} For every parent whose children are all filed, the
filed parent is compared to the signed sum of its filed children:
\begin{equation}
  \Big\lvert \text{parent} - \!\!\sum_{\text{children}}\!\! \text{sign}\cdot\text{child} \Big\rvert
  \;>\; \tau\cdot\lvert\text{parent}\rvert
  \quad\Rightarrow\quad \text{violation},
\end{equation}
with a tolerance of $\tau=5$ basis points of the parent magnitude.

\paragraph{Named identities.} Declared checks in \code{ref\_std\_identity\_check}
(each a left-hand item, a set of right-hand items with signs, a tolerance, and a
sector/standard scope) verify relations such as
$\text{Assets}=\text{Liabilities}+\text{Equity}$. When all operands resolve, the
signed residual is compared to the check's \code{tolerance\_bp}. Failures are
written to \code{ref\_std\_identity\_violation} with the LHS value, RHS value,
absolute delta, and delta in basis points --- an auditable record of where a
filer's statements do not tie out.

\begin{keybox}[Advisory, not blocking]
Identity violations are recorded for governance and surfaced to the committee's
data-quality report, but they do not by themselves discard data: real filings
occasionally fail an identity for benign reasons (rounding, disclosure
differences). The application's committee treats them as advisory by default
(see \emph{System Architecture \& Operations}); the warehouse keeps the evidence.
\end{keybox}

%======================================================================
\part*{Part IV\quad Metrics computation}
\addcontentsline{toc}{section}{Part IV\quad Metrics computation}

\section{The formula registry}
\label{sec:registry}
\emph{Code:} \code{metrics/formulas.py}, \code{metrics/compute.py}. Metrics are
defined declaratively as executable formulas over standardized line items,
generated from a line-item/metric registry so the public ids are governed in one
place. There are two derived layers:

\begin{itemize}
  \item \textbf{L1 line items} (\code{\_L1\_FORMULAS}) --- derived
        \emph{balance-sheet / income / cash-flow} quantities that behave like line
        items (e.g.\ \code{free\_cash\_flow}, \code{invested\_capital},
        \code{net\_operating\_assets}, \code{enterprise\_value},
        \code{tangible\_book\_equity}).
  \item \textbf{L2 metrics} (\code{\_METRIC\_FORMULAS}) --- ratios, margins, growth
        rates, turnovers, coverage, yields, and composite signals (e.g.\
        \code{return\_on\_invested\_capital}, \code{free\_cash\_flow\_yield},
        \code{price\_to\_earnings\_trailing}, \code{piotroski\_f\_score},
        \code{sloan\_accruals\_ratio}).
\end{itemize}

\begin{center}\small
\begin{tabular}{@{}>{\ttfamily\footnotesize\color{navy}}l >{\ttfamily\footnotesize}p{9.6cm}@{}}
\toprule
\normalfont\tabhead{Item / metric} & \normalfont\tabhead{Formula (registry)} \\
\midrule
free\_cash\_flow      & cash\_flow\_from\_operations $-$ |capital\_expenditures| \\
invested\_capital     & total\_equity $+$ net\_debt \\
net\_operating\_profit\_after\_tax & earnings\_before\_interest\_taxes $\times$ (1 $-$ cash\_tax\_rate) \\
return\_on\_invested\_capital & nopat / invested\_capital \\
gross\_margin         & gross\_profit / revenue \\
free\_cash\_flow\_yield & free\_cash\_flow / market\_capitalization \\
price\_to\_earnings\_trailing & market\_cap / net\_income\_attributable\_to\_common \\
days\_sales\_outstanding & accounts\_receivable\_net $\times$ 365 / revenue \\
\bottomrule
\end{tabular}
\end{center}

\subsection{Dependency ordering and safe evaluation}
Formulas reference other derived items (\code{return\_on\_invested\_capital} needs
\code{nopat} and \code{invested\_capital}, which need
\code{cash\_tax\_rate} and \code{net\_debt}). The registry is compiled into a
directed dependency graph and \emph{topologically ordered}; a cycle is a hard
error. Each formula is evaluated in a sandboxed namespace with a fixed set of safe
helpers --- \code{\_div} (null- and zero-safe division), \code{pct\_change},
\code{cagr}, \code{\_tax\_rate}, \code{\_abs\_capex} --- and no builtins, with a
missing operand resolving to \code{None} so a partial input yields \code{None}
rather than an exception. A result that is \code{NaN} or infinite is dropped.

\subsection{Point-in-time context and period joins}
Ratios that mix flows with market data or prior-period values are assembled with
care. Market capitalization is computed point-in-time as
$\text{price}(\text{period\_end})\times\text{shares}$; year-over-year and multi-year
formulas join the correct prior annual period ($\text{prev}$, $3y$, $5y$) via
\code{*\_prev} / \code{*\_3y} / \code{*\_5y} operands resolved against the entity's
own annual history. A handful of ``special'' metrics (momentum, volatility) come
from the price/market path rather than the fundamentals formulas. Each metric row
stores its \code{formula}, value, currency, \code{metric\_type}, category,
importance, unit type, and a \code{fallback\_applied} flag.

\section{International equities}
Non-US/JP equities follow a parallel, lighter path
(\code{metrics/compute\_intl.py}) that computes the same metric ids from
Yahoo-sourced fundamentals (\code{fact\_yahoo\_*}) into \code{fact\_metrics\_intl}.
It reuses the formula registry where the inputs are available and degrades where
they are not, so the screener and quant panels can include international names on
the same metric vocabulary.

%======================================================================
\part*{Part V\quad Reconciliation and validation}
\addcontentsline{toc}{section}{Part V\quad Reconciliation and validation}

\section{The reconciliation layer}
\label{sec:recon}
\emph{Code:} \code{metrics/recon.py} $\to$ \code{fact\_metrics\_recon\_\{us,jp\}}.
Reconciliation re-expresses every computed metric with its actual inputs and traces
each input back to the raw source. For a metric row it records:

\begin{center}\small
\begin{tabular}{@{}>{\ttfamily\footnotesize\color{navy}}l p{9.3cm}@{}}
\toprule
\normalfont\tabhead{Column} & \normalfont\tabhead{Content} \\
\midrule
formula              & The symbolic formula, e.g.\ \code{\_div(free\_cash\_flow, market\_capitalization)}. \\
formula\_with\_values & The same formula with each operand substituted by its value and name, e.g.\ \code{\_div(1.84e10 [free\_cash\_flow], 2.7e12 [market\_capitalization])}. \\
input\_values        & The resolved numeric value of every operand (\code{None} where missing). \\
source\_line\_items  & The standardized line items the operands resolved to. \\
source\_concept\_ids & The raw XBRL concepts behind those line items. \\
source\_filing\_ids  & The filings those facts came from. \\
raw\_trace           & Per-operand records: line item, value, source concept, filing, form, filed date, concept path, period role (current / prev / 3y / 5y), and (JP) context. \\
trace\_quality       & A tier describing how completely the metric traces to source (below). \\
\bottomrule
\end{tabular}
\end{center}

\subsection{Trace quality tiers}
Each reconciled metric is graded by how fully its inputs trace to filed source:

\begin{center}\small
\begin{tabular}{@{}>{\ttfamily\color{navy}}l p{9.8cm}@{}}
\toprule
\normalfont\tabhead{trace\_quality} & \normalfont\tabhead{Meaning} \\
\midrule
full\_source\_trace       & Every operand resolves to a filed source fact. \\
partial\_computed\_inputs & Some operands are themselves derived (L1) rather than filed; values are present. \\
partial\_missing\_inputs  & One or more operands are missing (the metric may still compute from the rest). \\
computed\_only            & Inputs exist as values but none resolve to a source fact (fully derived). \\
\bottomrule
\end{tabular}
\end{center}

To build the trace, recon reconstructs the same L1 line-item time series the
metrics stage used, resolves each formula token to its period-appropriate source
(current vs.\ \code{prev}/\code{3y}/\code{5y}), and de-duplicates the resulting
source records. The build is entity-chunked to bound memory, and has a freshness
gate so it is skipped when the recon table already covers the metrics table.

\begin{modelbox}{Worked example — free cash flow yield}
\small
For one entity-period, \code{free\_cash\_flow\_yield}$=$\code{\_div(free\_cash\_flow,
market\_capitalization)}. Recon records:
\begin{itemize}\setlength{\itemsep}{1pt}
  \item \code{formula\_with\_values}: \code{\_div(1.84e10 [free\_cash\_flow], 2.70e12 [market\_capitalization])}
  \item \code{free\_cash\_flow} is itself L1 (\code{cash\_flow\_from\_operations $-$ |capital\_expenditures|}); its operands trace to the filed \code{us-gaap:NetCashProvidedByUsedInOperatingActivities} and \code{us-gaap:PaymentsToAcquirePropertyPlantAndEquipment} concepts in a specific 10-K filing.
  \item \code{market\_capitalization} is computed point-in-time from price $\times$ shares.
  \item \code{trace\_quality}: \code{partial\_computed\_inputs} --- the fundamentals trace to source, but market cap is a derived price$\times$shares quantity, not a filed fact.
\end{itemize}
Every figure the committee cites can thus be unfolded to the filings and concepts
underneath it.
\end{modelbox}

\section{Validation and quality gates}
\label{sec:validate}
\emph{Code:} \code{quality/validate.py}. A compact validation summary is produced
per jurisdiction, spanning the whole pipeline:

\begin{itemize}
  \item \textbf{Layer counts} --- rows, entities, and distinct
        concepts/line-items/metrics at each of the raw, standardized, metric, and
        recon layers, so shrinkage or gaps between layers are visible at a glance.
  \item \textbf{Master quality} --- hard gates on the company master
        (\code{dim\_company\_\{us,jp\}}): e.g.\ US excludes ETFs/ETNs, funds/trusts,
        and non-operating SIC codes; JP enforces primary-ticker format, ISIN
        format, security-code/ticker consistency, and ticker-link integrity.
        A failed gate raises rather than silently shipping bad master data.
  \item \textbf{Source and XBRL state} --- per-source download/extract/parse
        progress and parse-error counts from \code{source\_filing\_state}, including
        US XBRL linkbase (cal/pre/def/lab) acquisition status.
  \item \textbf{Structural sparsity} --- coverage of the structural columns
        (statement type, parent/root, presentation order, weights) that graph
        closure depends on.
  \item \textbf{Recent runs} --- the last pipeline stage runs with row-in/row-out
        counts and errors, from \code{pipeline\_stage\_run}.
\end{itemize}

\section{How the standardized layer feeds the application}
The application reads only downstream of this pipeline. The committee's evidence
engine reads standardized fundamentals and metrics; the quant desk builds its panel
from \code{fact\_metrics\_*}; the WACC and factor models read the factor loadings
derived from prices. The committee also ships a read-only \emph{data-quality agent}
that scans exactly the five layers this document describes --- raw, standardized,
metrics, recon, and a cross-check against an independent source --- reports layer
scores and severity-ranked findings, and (in the mapping-check path) surfaces
raw concepts for an entity that went unmapped or resolved to an incompatible
sector, routing any proposed mapping change to a review queue rather than mutating
the governed table. The reconciliation trace of \S\ref{sec:recon} is what lets that
agent, and any analyst, unfold a headline number to the filings beneath it.

\vfill
\begin{center}\small\color{muted}
The standardized layer is the contract between heterogeneous filings and a single,
auditable analytical vocabulary. Every number the application shows can be traced
to the filing it came from, or is explicitly marked as derived.
\end{center}

\end{document}
